\documentclass[11pt]{article}
\usepackage{amsmath,amssymb,amsthm,mathtools}
\usepackage[margin=1in]{geometry}
\usepackage{hyperref}

\newtheorem{theorem}{Theorem}
\newtheorem{proposition}[theorem]{Proposition}
\theoremstyle{definition}
\newtheorem{remark}[theorem]{Remark}
\newtheorem{result}[theorem]{Computed Result}

\newcommand{\Z}{\mathbb{Z}}
\newcommand{\Q}{\mathbb{Q}}
\newcommand{\R}{\mathbb{R}}
\newcommand{\zetaK}{\zeta_{K}}
\newcommand{\Icos}{\mathcal{I}}
\newcommand{\fq}{\mathfrak{q}}

\title{The Weil Witness from the Icosian Triad:\\ the Gap, Localized}
\author{(VFD programme --- working draft)}
\date{2026}

\begin{document}
\maketitle

\begin{abstract}
The companion papers establish the icosian object exactly: the maximal icosian order
$\Icos$ on the $600$-cell over $K=\Q(\sqrt5)$, its triad $(\Icos, C_\varphi, N_H)$ with
Galois symmetry $\sigma$, and its Eisenstein theta $L$-function
$L(\Theta_\Icos,s)=\zetaK(s)\zetaK(s-1)$. Here we show that the same object also supplies
the prime-side ingredient of a Weil positivity witness $W(h)$, parameter-free, and we
locate precisely the single step it does not close. The prime residual is assembled
entirely from the geometric eigenvalues of the self-adjoint icosian Brandt operator at
level $\mathfrak p_{31}=(5\varphi-2)$ (the object's \emph{cuspidal} face), with no fitting
and no zero-data; the archimedean term is the cuspidal $L$-function's completion (gamma
factor and conductor), fixed by its functional equation. The implementation is validated
on $\zeta$ (where the zeros are known) and gives $W(h)\ge 0$ on the finite Gaussian test
family tried. \textbf{We make no claim about the Riemann Hypothesis.} Crucially: the
witness lives on the \emph{cuspidal} face, so its positivity quantifier is GRH for that
level-$31$ $L$-function, \emph{not} classical RH (classical $\zeta$ sits in the separate
Eisenstein face). The verified tests are moreover archimedean-dominated, so they are
consistent with, but not a sharp test of, that quantifier. The result is a
\emph{localization}, not a reduction of classical RH.
\end{abstract}

\section{The two arithmetic faces of one object}
The icosian order $\Icos$ carries two distinct arithmetic objects, and keeping them
separate is the whole point of this note.

\paragraph{Eisenstein face (theta).} The theta series of $\Icos$ has
$L(\Theta_\Icos,s)=\zetaK(s)\zetaK(s-1)$ exactly \cite{IcosianTriad}; via
$\zetaK=\zeta\cdot L(s,\chi_5)$ the Riemann $\zeta$ appears as one of its four factors.
This face \emph{contains} classical $\zeta$. Its parameters are $a_q=N(q)+1$, which
violate the Ramanujan bound $|a_q|\le 2\sqrt{N(q)}$, so it is \emph{not admissible as the
unitary cuspidal Satake input} used in the witness below. (It still has an explicit
formula of its own; we make no use of one here.)

\paragraph{Cuspidal face (Brandt).} The same order carries a self-adjoint Brandt/Hecke
operator $B(\fq)$ at level $\mathfrak p_{31}=(5\varphi-2)$, whose eigenvalues $a_\fq$ are
read directly from the icosian Brandt matrices---no point-counting, no fitting---and were
independently checked out-of-sample to equal the corresponding point counts
\cite{BrandtCuspidal}. The computed $a_\fq$ obey the Ramanujan bound in the tested range;
global Ramanujan for the associated form follows from its Hilbert-modular /
Jacquet--Langlands identification, not from the finite Brandt matrices alone. This
cuspidal face is where the witness lives.

\section{The Weil witness, and that its prime side is built from the object}
Fix the cuspidal level-$31$ $L$-function of the previous paragraph, with completed form
$\Lambda(s)=L_\infty(s)L(s)$ ($L_\infty$ its gamma factor, conductor $775$). For a test
function $h$ even and of \emph{positive type} (the Fourier transform $\widehat h\ge0$),
Weil's explicit formula reads, schematically,
\[
   W(h)\;=\;\underbrace{\mathrm{ARCH}(h)}_{\text{from }L_\infty,\ \text{conductor}}
   \;-\;\underbrace{\mathrm{PRIME}(h)}_{\text{from the }a_\fq}
   \;=\;\sum_{\gamma} h(\gamma),
\]
the sum over the nontrivial zeros $\rho=\tfrac12+i\gamma$ of $L(s)$ \cite{Weil}. The
witness $W$ is unconditional; Weil's criterion is the theorem
\[
   \text{GRH for this cuspidal } L \;\Longleftrightarrow\; W(h)\ge 0
   \quad\text{for every } h \text{ of positive type.}
\]

\begin{proposition}[the prime side is parameter-free from the object]\label{prop:witness}
In the computation \cite{BrandtCuspidal}, $\mathrm{PRIME}(h)$ is assembled entirely from
the geometric Brandt eigenvalues $a_\fq$ of $\Icos$'s cuspidal face, over the cached
prime ideals $N(\fq)\le 150$; the archimedean term $\mathrm{ARCH}(h)$ is the cuspidal
$L$-function's completion (gamma type and conductor $775$), fixed by its functional
equation. No constant is fitted, and no zero-data of the target enters the construction.
\end{proposition}

So the \emph{prime side} of the witness is derived from the icosian triad---it is not
placed beside it. The archimedean side is the automorphic completion of the chosen
cuspidal form, not itself read off the geometry; we state this division of inputs
explicitly rather than claiming the whole witness is geometric.

\section{What is computed, and the control}
The construction (\texttt{weil\_wall.py}, sourcing \texttt{geometric\_aP.py}) was run as
follows. All tests are Gaussian, $h(r)=e^{-r^2/2\sigma_h^2}$ with $\sigma_h\in\{3,4,5,6\}$.

\begin{result}[validation on $\zeta$]\label{res:gate}
Run on the Riemann $\zeta$, where the zeros are known, the same implementation reproduces
$\sum_\gamma h(\gamma)$ to relative error $\sim10^{-2}$ down to $10^{-6}$ across
$\sigma_h=3,\dots,6$. This validates the normalizations and pole/gamma bookkeeping; no
constants are fitted. (The $\zeta$ run uses $\zeta$'s own prime term, pole, and gamma
factor; it shares the framework, not the data.)
\end{result}

\begin{result}[positivity on the tested family]\label{res:pos}
For the cuspidal level-$31$ $L$-function, $W(h)\ge 0$ on every test in the family
($\sigma_h=3,4,5,6$; $W=+1.4047,+3.6961,+6.3931,+9.4125$).
\end{result}

\begin{result}[control]\label{res:ctrl}
Feeding the \emph{Eisenstein} bookkeeping $a_q=N(q)+1$ (the theta-face parameters) into
the same residual yields $32$ Ramanujan violations: these coefficients are not the
admissible cuspidal Satake input the residual requires. The control thus has teeth---the
witness is not satisfied by arbitrary input---without implying the Eisenstein face is
disqualified from all explicit-formula positivity.
\end{result}

\section{The gap, named precisely}
Two statements fix the honest scope.

\begin{remark}[the quantifier is GRH for the cuspidal $L$, not classical RH]
$W(h)\ge 0$ for \emph{all} $h$ of positive type is, by Weil's criterion, GRH for this
cuspidal level-$31$ $L$-function---\emph{not} classical RH. Classical $\zeta$ lives in the
Eisenstein face (\S1), which is not the witness face. The verified positivity is on a
finite Gaussian family; the universal quantifier is the wall and is GRH-equivalent, not
closable without the (out-of-reach) operator/positivity proof.
\end{remark}

\begin{remark}[the tests are archimedean-dominated]
Against our own result: in the tested family the prime residual is tiny relative to the
archimedean term (ratio $\sim10^{-4}$ down to $10^{-16}$ as $\sigma_h$ grows), and the
prime term is truncated to the cached range $N(\fq)\le150$ with no tail bound. Hence
$W(h)\ge0$ here is \emph{consistent with} GRH for this $L$-function but is \emph{not a
sharp test} of the positivity quantifier; we supply no sharp test family.
\end{remark}

\section{Conclusion}
The icosian triad is exact at the object level, and it supplies---parameter-free, from
its cuspidal Brandt face---the prime side of a Weil positivity witness $W(h)$ for a
specific cuspidal level-$31$ $L$-function, with the archimedean side taken from that
form's completion. The implementation is validated on $\zeta$ and gives $W(h)\ge0$ on
every test in a finite Gaussian family. The step not closed is the universal positivity
quantifier, which is \textbf{GRH for that cuspidal $L$-function}; and even the verified
positivity sits in an archimedean-dominated, finitely-truncated regime that is not a
sharp test of it.

The localization is therefore \emph{not} a reduction of classical RH. After choosing the
cuspidal Brandt face, the remaining wall is GRH for that one $L$-function---one named
inequality on one parameter-free prime residual. Classical $\zeta$ remains in the
separate Eisenstein face, which is not positivity-valid for this residual. We make no
claim about the Riemann Hypothesis. One object, two arithmetic faces---Eisenstein (theta,
containing $\zeta$) and cuspidal (Brandt, carrying the witness)---and the wall is the
positivity quantifier on the second.

\begin{thebibliography}{9}
\bibitem{IcosianTriad} VFD programme, \emph{The Icosian Triad on $V_{600}$}
  (\texttt{icosian-triad-v600} release, v1.1): $L(\Theta_\Icos,s)=\zetaK(s)\zetaK(s-1)$
  exactly, verified in the maximal order; states it does not address RH/GRH.
\bibitem{BrandtCuspidal} VFD programme, route-B icosian Brandt computation
  (\texttt{icosian\_brandt\_cuspidal\_geometry/route\_b}, \texttt{weil\_wall.py},
  \texttt{geometric\_aP.py}): geometric Hecke eigenvalues $a_\fq$ of the self-adjoint
  icosian Brandt operator at level $\mathfrak p_{31}$, cache $N(\fq)\le150$,
  out-of-sample checked against point counts; the witness run uses conductor $775$ and is
  archimedean-dominated.
\bibitem{Weil} A.~Weil, \emph{Sur les ``formules explicites'' de la th\'eorie des
  nombres premiers}.
\end{thebibliography}

\end{document}
